\hypersetup{pageanchor=false}
\pagenumbering{gobble}
\pagestyle{empty}

\chapter*{Declaration}
\thispagestyle{empty}
I hereby declare that this thesis is my own work and has not been submitted in any form for another degree or qualification at any university or institution.

\clearpage
\chapter*{Acknowledgements}
\thispagestyle{empty}
I would like to thank Dr. Christodoulos Stylianou, my supervisor, for guiding, giving constructive feedback and supporting me in the development of this project. His guidance enabled me to focus the work, enhance the technical aspects of implementation, and structure my thesis. Finally, I want to express my gratitude to the Department of Computer Science and Engineering at European University Cyprus for the academic environment and the resources it offered in the course of this work.

I would also like to thank the Department of Computer Science and Engineering at European University Cyprus for providing the academic environment and resources that supported this work. I am grateful to the lecturers and staff who contributed to my studies and helped me build the knowledge needed to complete this project.

Finally, I would like to thank my family and friends for their encouragement, patience, and support during the preparation of this thesis. Their motivation and understanding were important throughout the research, implementation, and writing process.

\clearpage
\chapter*{Abstract}
\thispagestyle{empty}
\acp{LLM} are now playing a crucial role in various applications of \ac{NLP}, such as question answering, summarization, and information retrieval. However, the high computational and operational cost of inference, document retrieval, indexing and data handling should be taken into account when designing a system. One way to address this issue is by combining a language model with an external source of knowledge, which can then be used to retrieve relevant documents and “feeding” them to the language model as a context before generating an answer: \ac{RAG}. 

This thesis explores deployment and evaluation of \ac{RAG} pipelines, particularly in terms of retrieval-level optimization, reproducibility and deployment portability. The experimental design is staged. The first phase consists of six pipelines implemented and evaluated locally on a personal computer: an in-memory pipeline as a baseline, a Milvus vector-database pipeline with an additional tuned one, a dense retrieval pipeline with reranking, a hybrid retrieval pipeline, and a hybrid retrieval pipeline with reranking. In the second stage, the same implementations, supporting files, scenario queries and result-logging structure will be run on a High-Performance Computing (HPC) environment. This enables the comparison between the variants of the RAG in each of them. 

The evaluation framework takes into account the performance of the system as well as the quality of the output. The system level measurements are the following: \textit{index time}, \textit{retrieval time}, \textit{generation time}, and \textit{total execution time}. The evaluation of output-quality takes into account aspects such as the relevance of the retrieved documents, the correctness of the answer, the groundedness, the completeness, and the risk of hallucination. This is relevant as the speed with which a RAG system returns information is not necessarily its most important characteristic, it should also be able to provide useful evidence to enable the generation of responses which are backed by evidence. The thesis provides a structured local comparison of the implemented RAG variants, it provides a reproducible workflow including a run and log experiment, and it provides a framework for extending the same files to running experiments on HPC. Local deployment and comparison are the baseline at this stage. The next experimental stage is the implementation of the HPC and a comparison between local and HPC; this will be completed in the results chapter with the appropriate measurements. Full model training is not included in the main, as is full model fine-tuning and large-scale multi-node inference experimental scope.

\clearpage
\chapter*{List of Abbreviations}
\thispagestyle{empty}
\begin{acronym}[BERTScore]
\acro{NLP}{Natural Language Processing}
\acro{LLM}{Large Language Model}
\acro{RAG}{Retrieval-Augmented Generation}
\acro{HPC}{High-Performance Computing}
\acro{GPU}{Graphics Processing Unit}
\acro{CPU}{Central Processing Unit}
\acro{KV}{Key--Value}
\acro{DPR}{Dense Passage Retrieval}
\acro{BM25}{Best Matching 25}
\acro{HNSW}{Hierarchical Navigable Small World}
\acro{FAISS}{Facebook AI Similarity Search}
\acro{TP}{Tensor Parallelism}
\acro{PP}{Pipeline Parallelism}
\acro{DP}{Data Parallelism}
\acro{FSDP}{Fully Sharded Data Parallel}
\acro{TGI}{Text Generation Inference}
\acro{TIS}{Triton Inference Server}
\acro{MMLU}{Massive Multitask Language Understanding}
\acro{HELM}{Holistic Evaluation of Language Models}
\acro{MLPerf}{Machine Learning Performance Benchmark Suite}
\acro{PFS}{Parallel File System}
\acro{GPFS}{General Parallel File System}
\acro{RAGAS}{Retrieval Augmented Generation Assessment}
\acro{ROUGE-L}{Recall-Oriented Understudy for Gisting Evaluation - Longest Common Subsequence}
\acro{BERTScore}{Bidirectional Encoder Representations from Transformers Score}
\acro{AWS}{Amazon Web Services}
\acro{I/O}{Input/Output}
\acro{vLLM}{Virtual Large Language Model}
\acro{FP32}{Single-Precision Floating-Point (32-bit)}
\acro{FP16}{Half-Precision Floating-Point (16-bit)}
\acro{INT4}{4-bit Integer Precision}
\acro{INT8}{8-bit Integer Precision}
\acro{GPTQ}{Generalized Post-Training Quantization}
\acro{AWQ}{Activation-aware Weight Quantization}
\acro{KV-cache}{Key-Value Cache}
\acro{KILT}{Knowledge Intensive Language Tasks}
\end{acronym}

\clearpage
\pagenumbering{arabic}
\setcounter{page}{1}
\hypersetup{pageanchor=true}
\pagestyle{frontmatter}

{\setstretch{1.0}
\begingroup

\titleformat{name=\chapter,numberless}
  {\raggedright\bfseries\fontsize{22}{26}\selectfont}
  {}
  {0pt}
  {}

\phantomsection
\markboth{Contents}{}
\tableofcontents

\clearpage
\phantomsection
\addcontentsline{toc}{chapter}{List of Figures}
\markboth{List of Figures}{}
\listoffigures

\clearpage
\phantomsection
\addcontentsline{toc}{chapter}{List of Tables}
\markboth{List of Tables}{}
\listoftables

\endgroup
}

\clearpage
\pagestyle{fancy}
